% Fichier complété par tools/complete_missing_corrections.py.
% Source : COR/COR-03-PI/65_Eclipse_02/65_Eclipse_02.tex
% Statut : rédigé (9 question(s)).
\begin{corrigebox}[Corrigé rédigé]
\CorrigeQuestion{1}
Le correcteur intégral ajoute un pôle à l'origine; le gain de boucle devient
            infini à basse fréquence et l'erreur statique d'échelon est nulle.
\CorrigeQuestion{2}
Après réduction au second ordre dominant, le temps optimal est obtenu pour
            $\xi\simeq0,69$, soit $t_{5\%}\simeq3/(\xi\omega_0)$. On détermine $K_i$ en
            identifiant le polynôme caractéristique à cette forme; la valeur numérique est
            celle qui place les pôles dominants sur cette condition.
\CorrigeQuestion{3}
Le Bode corrigé est celui du procédé auquel on ajoute $20\log K_i$ et
            $-20$ dB/décade, avec une phase supplémentaire de $-90^\circ$ due à $1/p$.
\CorrigeQuestion{4}
La marge de phase est $180^\circ+\arg L(j\omega_c)$ à la pulsation où le
            gain vaut 0 dB.
\CorrigeQuestion{5}
La valeur limite de $K_i$ translate le gain jusqu'à ce que la coupure se
            produise à la pulsation où la phase vaut $-180^\circ+M_{\varphi,min}$. Ainsi
            $K_{i,lim}=1/|L_0(j\omega_{lim})|$.
\CorrigeQuestion{6}
Sur le diagramme donné pour $K_i=7000$, on lit le gain à la pulsation de
            phase cible puis on corrige proportionnellement :
            $K_{i,lim}=7000\,10^{-G_{dB}(\omega_{lim})/20}$.
\CorrigeQuestion{7}
Une valeur beaucoup plus élevée que l'optimum de rapidité augmente la
            bande passante, le dépassement et la sensibilité au bruit; elle n'est donc pas
            retenue même si la stabilité limite est respectée.
\CorrigeQuestion{8}
Un double intégrateur rendrait nulles les erreurs d'échelon et de rampe,
            donc améliorerait la précision basse fréquence.
\CorrigeQuestion{9}
Mais il ajoute $-180^\circ$ de phase : sans zéros compensateurs, la marge
            devient négative ou très faible. Il ne permet donc pas d'assurer la stabilité.
\end{corrigebox}
